%% =====================================================================
%%  T-10-05  The Withdrawal
%%  -------------------------------------------------------------------
%%  One idea per file. Core is mandatory. Slots in canonical order.
%%  Max 700 words. No layout commands. No filler. New source material
%%  for this entry goes in sources/B3/T-10-05.tex -- never by
%%  rewriting this file (docs/RULES.md R0.1).
%%  =====================================================================
%% @kind: topic
%% @id: T-10-05
%% @title: The Withdrawal
%% @subtitle: Use absence --- presence taken for granted is presence priced at zero
%% @chapter: c10
%% @order: 50
%% @status: stable
%% @sources: B3
%% @updated: 2026-09-19

\topic{T-10-05}{The Withdrawal}{Use absence --- presence taken for granted is presence priced at zero}

\begin{core}
Too much presence makes a person common: the constantly available are taken for granted, and what is taken for granted is not respected. The move is the engineered withdrawal --- starve the audience of your presence for a measured interval, create the pattern of presence and absence, and let scarcity re-price what habit had discounted to nothing.
\end{core}
\begin{mechanism}
B3's case is the judge who served everyone so faithfully that they grew to take his services for granted; only by withdrawing completely --- letting the Medes taste life without him --- did he regain the veneration, and he returned to a throne built on the lesson. The mechanism is habituation: value is set at the margin, and a person who is always there has no margin. The source extends the law of scarcity to skills --- make what you offer rare and hard to find, and its value instantly rises --- and adds the colder variant: the threat of absence, the possibility that you might be lost for good, forces respect without a single day's withdrawal. The reversal is the guard-rail: absence only works after presence has been established. Leave too early and you are not missed; you are simply forgotten. Withdrawal is the accent, not the foundation.
\end{mechanism}
\begin{conditions}
Strongest where you are already valued but increasingly expected --- the reliable supplier, the ever-available friend, the fixture. Strongest also in attention economies where novelty is currency. It fails at the entry stage (unknowns withdraw into unknowns), in contracted obligations (absence there is breach, not strategy), and against counterparts who happily fill your absence with a substitute --- audit substitutability first (T-11-01's logic applies to its author too).
\end{conditions}
\begin{application}
  \item Diagnose honestly: are you being valued, or merely consumed? Taken-for-granted has a texture --- requests arrive without acknowledgment.
  \item Withdraw on a pattern, not in a huff: shorter responses, fewer appearances, a named unavailability. Sulks re-price you as volatile, not scarce.
  \item Let the absence be felt where it matters --- one skipped ritual outperforms a dozen small absences.
  \item Return with something new. The withdrawal re-opens attention; the return must spend it on a re-introduction.
  \item With skills, withdraw by rationing: the scarce offer, not the absent one.
\end{application}
\begin{feedback}
  \item Working: the messages change tone --- from instructions to invitations.
  \item Working: your price moves, or the terms improve, without you raising either.
  \item Failing: the absence passes unnoticed. The re-pricing has revealed a market price of zero; now you know (T-11-01).
  \item Failing: the counterpart retaliates by out-withdrawing you. The pattern became a contest of sulks.
\end{feedback}
\begin{failure}
The tactic corrupts into game-playing --- relationships run on engineered absence become chess, and the counterpart who finally notices the pattern reads everything afterwards as manipulation, including the genuine parts. It also carries real cost: absences are not free; a person can lose standing, information flow (T-27-03's warning), and momentum during a withdrawal that outlasts the audience's memory.
\end{failure}
\begin{countermeasures}
  \item Notice the engineered absence being run on you: the person whose availability dropped just after you stopped acknowledging them. Decide on merits whether to chase.
  \item Distinguish scarcity from dysfunction: a colleague who vanishes to re-price themselves and one who vanishes from dysfunction look identical for a month. Judge by the return.
  \item In your own house, watch the always-available. The price of reliable presence is invisibility; compensate your fixtures deliberately, in words, before the market does it in silence.
  \item If someone's withdrawal is a sulk in strategy's clothing, do not bid. Bidding a sulker back sets the price for every future sulk (T-19-02).
\end{countermeasures}

\sources{B3}
\seesources
\seealso{T-10-01, T-10-04, T-27-03}
